\hypertarget{schedules-and-initial-classification}{%
\section{Schedules and initial classification}\label{method-financing-methods-schedules-and-initial-classification}}

The following liability and ownership methods use selected dated AASB 9 and AASB 132 provisions \citep{aasb9dated,aasb132}. Issuer, holder, contractual cash flow and accounting measurement remain separate throughout the schedule.

DebtContract contains parties, currency, principal, issue/settlement dates, payment calendar, fixed/index/spread/reset terms, compounding/day-count, business-day adjustment, principal changes, fees, options, security, priority and legal default/remedy terms. Each term has evidence/admission provenance. The economic schedule emits due cash independently of the recognition method. A market-value adapter may value that schedule but cannot choose amortised cost, FVOCI, FVTPL or liability/equity classification.

An investor's classification evaluates business model and contractual cash-flow features; a debt label alone cannot establish the AASB 9.4.1.2/4.1.2A conditions. At initial recognition the selected 5.1.1 rule measures fair value and qualifying costs as appropriate. An unknown day-one observable-value condition yields unresolved under 5.1.1A/B5.1.2A. The retained acquisition cases distinguish FVTPL expensed fees, amortised-cost fees, original EIR and POCI credit-adjusted EIR. No rate method inferred from market value can overwrite the contract rate.

For an issuer with a plain amortised-cost debt, net proceeds 100 and one legally due payment 121 at the end of two annual periods, the exact EIR solves 100=121/(1+r)\^{}2, hence r=10\%. Issue debit cash 100/credit debt 100; year-one debit interest 10/credit debt 10; year-two debit interest 11/credit debt 11; redemption debit debt 121/credit cash 121. Outstanding contractual payment stays 121 until settled even while carrying value is 100 or 110. Total interest 21 equals cash outflow 121 minus proceeds 100. An immediate transaction fee 2 changes net proceeds to 98 and the EIR root; retaining 10\% is wrong. The numerical root method is in the numerical-method discussion; an irrational rate is not silently an exact rational.

A fractional period uses the declared convention. Annual-effective 21\% over half a year produces factor 1.1 and interest 10 on 100; simple proration gives 10.5 and is wrong for that contract. A floating coupon resets from the admitted fixing available at the reset date, not a later observed rate. Actual day-count and holiday-calendar versions are immutable inputs. Unknown fixings preserve conditional paths, rather than defaulting to the latest value.

\hypertarget{modification-extinguishment-and-default}{%
\section{Modification, extinguishment and default}\label{method-financing-methods-modification-extinguishment-and-default}}

At a modification, retain original and replacement terms, counterparty, fees, effective date, legal-release evidence and original EIR. Determine derecognition before measurement. AASB 9.3.3.1--3.3.3 and B3.3.1--B3.3.6 govern the selected issuer-liability cases. The quantitative ten-percent test compares the PV of new terms, including specified borrower/lender fees, at original EIR against remaining original PV. It is not a universal asset test, nor proof that every change below ten percent is immaterial; changed currency, rights and other qualitative terms need the applicable analysis. Insufficient evidence remains unresolved. Third-party payment without legal release does not extinguish debt.

On the first anniversary above, old remaining 121 has PV 110. Replace with 133.1 due at the original maturity date, one year after the modification, with no fees and an explicit supported substantial-terms conclusion: PV 121 differs by exactly 10\%, satisfying the inclusive test. Assume separately supported new fair value 118; debit old debt 110, debit extinguishment loss 8, credit new debt 118. The new EIR is 133.1/118\ensuremath{-}1, not automatically the old rate. If the replacement is instead a stated non-substantial same-currency reduction to 115.5, no fees or other changes, PV 105 leads to debit debt 5/credit modification gain 5 under the B5.4.6 recalculation mechanism, then interest 10.5 and settlement 115.5. The actual change below ten percent still needs the qualitative finding. Repurchase of carrying amount 110 for 108 extinguishes the repurchased obligation: debit debt 110/credit cash 108/credit gain 2. Unamortised fees belong in the carrying amount before that comparison, not a second gain adjustment.

If cash available is 80 against debt 121 due, record the due obligation and an unpaid/default economic state; payment 80 leaves contractual unpaid 41 only when that partial payment actually occurs. Do not cancel the remaining liability, recognize a waiver or issue a loan unless a separately admitted event authorizes it. Covenant acceleration, cross-default, cure periods, collateral realization and restructuring require the actual legal terms and priority waterfall. A scenario may explore specified alternatives with explicit assumptions; it may not silently select the cheapest settlement as an accounting fact.

\hypertarget{equity-and-compound-instruments}{%
\section{Equity and compound instruments}\label{method-financing-methods-equity-and-compound-instruments}}

EquityInstrument retains class, voting/distribution/liquidation rights, issued, treasury and outstanding counts, par/legal-capital requirements and conversion terms. Share counts are integer quantities separate from currency balances. Issue 100 ordinary shares at 2 with no fees: debit cash 200/credit equity 200, issued and outstanding 100. Repurchase 10 at 3 for treasury: debit treasury-equity 30/ credit cash 30; outstanding 90, issued 100, treasury 10. Reissue 5 at 4: debit cash 20/ credit treasury-equity 15/credit other equity 5; outstanding 95, treasury 5. There is no trading profit under the selected AASB 132.33 rule; a statutory par split requires that jurisdiction's facts. Dividend declaration creates a liability only at the admitted point at which the obligation exists; proposal of a dividend is not itself payment or automatically an expense.

For a qualifying fixed-for-fixed issuer conversion option in functional currency, AASB 132.28--32 separates liability fair value and residual equity. Assume total issue proceeds 100, standalone debt fair value 90, zero coupon, maturity two years and cash redemption 100 or fixed 10 ordinary shares; no extra derivative, tax or fee. Issue debit cash 100/credit debt 90/credit conversion-equity 10. Effective interest accretes debt 90 to 100. At maturity conversion under original terms debits debt 100 and credits issued-share equity 100; the original equity 10 stays in equity (a permitted internal transfer is not income). Shares increase 10 and no redemption cash leaves. Redemption instead debits debt 100/credit cash 100; unexercised equity is not remeasured into profit. AG 32 supports no conversion gain/loss in that stated case. Variable share count, cash alternatives, foreign currency and changed conversion terms require classification/modification analysis; they cannot be forced into this example.

Financial statement mappings separate interest, modification/extinguishment results, owner transactions and cash settlement. A financing-liability note reconciles opening carrying amount, cash proceeds/repayments, effective interest, fees, conversions, modifications and FX. Total shares and earnings-per-share denominators need dates and dilution conditions; outstanding shares at year end are not a substitute for a weighted period denominator. Full option valuation, hedge accounting, EPS and legal-capital regimes remain explicitly additional methods; the proposed records preserve their required inputs.
